\newcommand\UniversiteAdi{Niğde Ömer Halisdemir Üniversitesi}
\newcommand\BolumAdi{MEKATRONİK BÖLÜMÜ}
\newcommand\DersKodu{MKT2002}
\newcommand\DersAdi{BİLGİSAYARLI KONTROL SİSTEMLERİ}
\newcommand\SinavAdi{Ödev 3}
\newcommand\SinavTarihi{10.03.2025}
\newcommand\SinavSaati{10:00}
\newcommand\SinavSuresi{90dk}
% Bölüm 6den sonra
\pagestyle{fancy}
\fancyhf{} % clear existing header/footer entries
\fancyhead[R]{Öğrenci No:\hspace{4.5cm}}
\fancyhead[L]{Ad Soyad:\hspace{7cm}}
\noindent \includegraphics[width=0.1\textwidth]{logo}
\begin{tabular}{
		p{0.15\linewidth}
		p{0.15\linewidth}
		p{0.2\linewidth}
		p{0.1\linewidth}
		p{0.15\linewidth}}
	\multicolumn{5}{c}{\textbf{\BolumAdi}}        \\
	\multicolumn{5}{c}{\textbf{\DersAdi}}         \\\hline
	\multicolumn{1}{|r|}{Ders Kodu:}  &
	\multicolumn{1}{|c|}{\DersKodu}   &
	\multicolumn{1}{|c|}{}            &
	\multicolumn{1}{|r|}{Tarih:}      &
	\multicolumn{1}{|c|}{\SinavTarihi}            \\\hline
	\multicolumn{1}{|r|}{Sınav Türü:} &
	\multicolumn{1}{|c|}{\SinavAdi}   &
	\multicolumn{1}{|c|}{}            &
	\multicolumn{1}{|r|}{Saat:}       &
	\multicolumn{1}{|c|}{\SinavSaati}             \\\hline
	\multicolumn{1}{|r|}{Dönemi:}     &
	\multicolumn{1}{|c|}{2024-2025}   &
	\multicolumn{1}{|c|}{}            &
	\multicolumn{1}{|r|}{Süre:}       &
	\multicolumn{1}{|c|}{\SinavSuresi}            \\\hline
	                                  &   &  &  & \\
\end{tabular}\\\\
\noindent\begin{center}
	\begin{tabular}{|r|c|}\hline
		               & \textbf{Toplam} \\\hline
		\textbf{Puan:} & \textbf{100}    \\\hline
		\textbf{Not:}  &                 \\\hline
	\end{tabular}\end{center}
\noindent\textbf{Uyarı:}
\begin{itemize}\bfseries
	\item Soruları dikkatlice okuyunuz. Hesap makinesi kullanılabilir.
	\item İşlemleri atlamadan ve ayrıntılı olarak veriniz. Sadece nümerik yanıtlar veya çizimler ara işlemler olmadan kabul edilmemektedir.
\end{itemize}
\noindent\textbf{Soru:} Blok diyagramı

\begin{figure}[!h]
	\centering
	\begin{tikzpicture}
		\node at (0, 0.5) (r) {$r$};
		\node at (0.8, 0.4) (plus) {$+$};
		\node at (1.1, -0.7) (minus) {$-$};
		\node at (3.8, 0) (fs) {$F(s)$};
		\node at (7.3, 0) (gs) {$G(s)$};
		\node at (9.3, 0.4) (y) {$y$};
		\node at (5.5, 0.4) (u) {$u$};
		\node at (2, 0.4) (e) {$e$};

		\draw [->](0,0) -- (1,0);
		\draw (1.4,0) circle (0.4);
		\draw [->](1.8,0) -- (2.8,0);
		\draw (2.8,-1) rectangle (4.8,1);
		\draw [->](4.8,0) -- (6.3,0);
		\draw (6.3,-1) rectangle (8.3,1);
		\draw [->](8.3,0) -- (9.3,0);

		\draw (8.75,0) -- (8.75,-2);
		\draw (8.75,-2) -- (1.4,-2);
		\draw [->](1.4,-2) -- (1.4,-0.4);

	\end{tikzpicture}
	\caption{P-tipi Birim Geribeslemeli Kontrol Yapısı}
	\label{fig:PTipiKontrolBlokDiyagrami}
\end{figure}
olarak verilmiştir.

\noindent\textbf{a)} $F(s)=k$ için kapalı çevrim transfer fonksiyonu
\begin{equation*}
	T(s)\triangleq \frac{y(s)}{r(s)}
\end{equation*}
olmak üzere $T(s)$ transfer fonksiyonunun matematiksel ifadesini elde ediniz.
